\documentclass[11pt]{article}
\usepackage{amsmath,amssymb,amsthm}
\usepackage[margin=1in]{geometry}

\newtheorem{theorem}{Theorem}
\newtheorem{lemma}{Lemma}
\newtheorem{proposition}{Proposition}

\begin{document}

\section*{Manuscript Insert: Cap-Released \(C^*\) Annulus Certificate}

The following text is intended for Paper 2 after author review.  It is a
standalone draft and does not modify the paper source.

This is a hybrid computer-assisted certificate.  The proof text below records
the mathematical reduction, the KKT and Hessian margins, the cutoff-zero
structure, and the rounding model.  The only external data required are the
archived cap-\(1000\) candidate, released-face Hessian, theorem-gate ledger,
and verification scripts.  The exploratory rank-window scans are not part of
the certificate.

For a shell-scale vector \(z\) in the finite face below, write \(v_z\) for the
corresponding divergence-free Fourier field and set
\[
   \Phi(z):=\frac{|B(v_z,v_z,\Delta v_z)|}{X(v_z)^2D(v_z)}.
\]

\begin{theorem}[finite one-high annulus certificate]
\label{thm:cstar-annulus-cap1000}
Let \(\mathcal F_{1413}\) denote the finite one-high annulus shell-split face
generated by the certified direction cache through the last active shell
\(|k|^2=2260\).  For the finite-face objective \(\Phi\),
\[
   \sup_{z\in \mathcal F_{1413}} \Phi(z) \le 0.31226427 .
\]
Moreover the certified candidate \(z_*\) satisfies
\[
   \Phi(z_*) = 0.31226426322102058 .
\]
\end{theorem}

\begin{proof}
The certificate is carried by three finite ledgers: the cap-released optimizer
output, the direct Hessian replay on the released free face, and the theorem
gate ledger.  The optimizer was rerun with scale cap \(1000\).  Its maximum
coordinate is \(200.00000000868158\), so no coordinate is at the cap.  The
artificial cap present in the earlier scale-\(200\) certificate is therefore
not part of the active proof.

The KKT inequalities are separated into lower-bound, upper-bound, and free
coordinates.  There are \(154\) lower-bound coordinates, and the largest
lower-bound gradient is
\[
   -3.5345695459004496\cdot 10^{-13}<0.
\]
There are no upper-bound coordinates.  On the \(647\) free coordinates,
\[
   \|\nabla_{\mathrm{free}}\Phi(z_*)\|_\infty
      \le 3.963733163849133\cdot 10^{-12}.
\]

The active Hessian was recomputed directly at the cap-\(1000\) point using the
alias-safe CUDA FFT replay on all \(647\) free coordinates.  Six of these
coordinates have zero Hessian diagonal.  They are not flat active directions:
the cutoff-zero audit verifies that they have zero one-high triads and zero
shared one-high triads, with total linear contribution below
\[
   5.5965044354606178\cdot 10^{-49}.
\]
Thus these six rows are removed by the structural no-triad lemma.

On the remaining \(641\)-dimensional nonzero block, the symmetrized Hessian is
strictly negative.  At diagonal threshold \(10^{-30}\), the largest eigenvalue
is
\[
   -3.0235227279516019\cdot 10^{-23}.
\]
Equivalently, after curvature normalization, the negative Hessian has minimum
eigenvalue
\[
   0.69033821039725485
\]
and Gershgorin row margin
\[
   0.50538223782086966.
\]
This gives a stable Schur/KKT certificate for the free face.

The residual quadratic gain allowed by the finite-gradient defect is far below
the display margin.  The candidate is below \(0.31226427\) by
\[
   6.778979433352816\cdot 10^{-9},
\]
while the tight residual-plus-target row in the theorem ledger leaves
\[
   5.9791946083210946\cdot 10^{-9}
\]
after absorbing the active quadratic residual.  Standard floating-point
rounding is accounted for with
\[
   \gamma_m=\frac{mu}{1-mu},\qquad u=2^{-53}.
\]
The denominator replay uses \(4838\) charged operations, and the remaining
post-residual numerator budget allows \(172463812\) charged operations before
the displayed target is exhausted.  The finite replay therefore lies below
\(0.31226427\) with the stated rounding model, completing the finite-face
certificate.
\end{proof}

\begin{proposition}[interface with the far tail]
\label{prop:cstar-annulus-tail-interface}
If the far-tail exclusion beyond \(|k|^2=2260\) is attached to
Theorem~\ref{thm:cstar-annulus-cap1000}, then the same displayed value
\[
   0.31226427
\]
is the corresponding full one-high annulus certificate.
\end{proposition}

\begin{proof}
The finite theorem closes all shell-split variables through the last shell in
the one-high table.  The only remaining directions are beyond
\(|k|^2=2260\).  The far-tail proposition excludes exactly those directions by
the separate no-first-variation and corrected-Hessian tail certificate.  Hence
no admissible one-high annulus perturbation can exceed the finite-face bound.
\end{proof}

\paragraph{Certificate artifacts.}
The finite theorem uses the cap-\(1000\) optimizer JSON, the cap-\(1000\)
released-face Hessian JSON, and the cap-\(1000\) theorem-gate JSON.  Full
filesystem paths and SHA256 hashes are listed in the accompanying proof packet.
\begin{itemize}
\item final candidate: \texttt{cap1000\_polish},
\item released-face Hessian: \texttt{cap1000\_free},
\item theorem-gate ledger: \texttt{cap1000\_verification}.
\end{itemize}

The recommended archive is intentionally minimal: it contains the three
certificate data files and the two verification scripts, not the exploratory
search history.

\end{document}
